% =============================================================================
% CHAPTER 8: STOCHASTIC SURPRISE GATE
% IP Traceability: IP-007 (Ai Surprise_small.pdf)
% =============================================================================

\chapter{Stochastic Surprise Gate}
\label{chap:stochastic-surprise}

\section{IP Document Summary}

\begin{tracerecord}
\textbf{IP Source ID} & \ipid{IP-007} \\
\textbf{Document} & Ai Surprise\_small.pdf \\
\textbf{Author} & Tom Pfaff, Craig Turrell \\
\textbf{Date} & April 2026 \\
\textbf{Pages} & 19 \\
\textbf{Abstract} & Stochastic Surprise Gate (SSG) framework for AI memory systems using Hull-White model dynamics \\
\end{tracerecord}

\subsection{Document Abstract}

\begin{ipsourcebox}{IP-007, Abstract}
\textit{``This paper proposes a novel theoretical framework that reframes memory management as a stochastic surprise detection and encoding problem. We draw a formal analogy between the surprise-driven memory gate in Google's Titan architecture and the shock-absorbing stochastic differential equations of the Hull-White financial model. This leads to our core contribution: the Stochastic Surprise Gate (SSG), a memory control mechanism in which the writing and retention strength for an item are proportional to its Z-scored deviation from a continuously updated sequence expectation.''}
\end{ipsourcebox}

\subsection{Key Contributions}

The IP document introduces:
\begin{enumerate}[nosep]
    \item \textbf{Stochastic Surprise Gate (SSG)}: Memory gating based on statistical significance
    \item \textbf{Hull-White Analogy}: Financial model applied to AI memory
    \item \textbf{Z-score as Surprise}: Normalized deviation as memory trigger
    \item \textbf{Audit Trail}: Complete statistical audit of memory decisions
\end{enumerate}

% =============================================================================
\section{Key Concepts Identified}
% =============================================================================

\begin{table}[H]
\centering
\caption{IP-007 Key Concepts and Implementation Targets}
\label{tab:ip007-concepts}
\begin{tabularx}{\textwidth}{@{}p{3.5cm}Xl@{}}
\toprule
\textbf{Concept} & \textbf{Description} & \textbf{Section} \\
\midrule
Z-score Surprise & Statistical significance as memory trigger & \ref{sec:ip007-zscore} \\
Audit Trail & Complete logging of statistical decisions & \ref{sec:ip007-audit} \\
Confidence Calibration & Historical correlation for trust metrics & \ref{sec:ip007-confidence} \\
Adaptive Calibration & Self-adjusting model parameters & \ref{sec:ip007-adaptive} \\
\bottomrule
\end{tabularx}
\end{table}

% =============================================================================
\section{Concept: Z-score as Surprise Metric}
\label{sec:ip007-zscore}
% =============================================================================

\begin{ipsourcebox}{IP-007, Innovation Section}
\textit{``The Stochastic Surprise Gate (SSG) represents an innovation at the intersection of quantitative finance and machine learning. By adapting the Hull-White interest rate model---specifically its stochastic shock term $\sigma dW(t)$---to neural memory systems, the SSG establishes a rigorous, probabilistic framework for memory gating. This formulation defines `surprise' as a statistically significant deviation from expectation, formalized through a Z-scored metric: $Z = |r_{obs} - E[r]| / \sigma$, where surprise $= \max(0, Z - \text{threshold})$.''}
\end{ipsourcebox}

\begin{tracerecord}
\textbf{IP Source ID} & IP-007 \\
\textbf{IP Location} & Innovation Section, Page 6 \\
\textbf{Concept} & Z-score as statistically rigorous surprise metric \\
\midrule
\textbf{Implementation Path} & \filepath{src/intelligence/analytics/anomaly\_detector.py} \\
\textbf{Key Files} & \filepath{anomaly\_detector.py} \\
\textbf{Functions/Classes} & \funcname{AnomalyDetector.detect\_single()} \\
\textbf{Adaptation Type} & Conceptual (memory gating $\rightarrow$ anomaly detection) \\
\end{tracerecord}

\begin{implbox}{src/intelligence/analytics/anomaly\_detector.py, Lines 162--222}
\begin{lstlisting}[language=Python]
def detect_single(
    self,
    value: float,
    reference_values: np.ndarray,
) -> AnomalyResult:
    """
    Detect if a single value is anomalous compared to reference distribution.
    
    Implements IP-007 Z-score surprise concept:
    Z = |r_obs - E[r]| / sigma
    
    Args:
        value: Single value to test (r_obs)
        reference_values: Reference distribution values
        
    Returns:
        AnomalyResult with Z-score (surprise metric)
    """
    # Determine method (Z-score vs MAD)
    method, _ = self._determine_method(reference_values)
    
    # Calculate Z-score (the "surprise" metric from IP-007)
    if method == DetectionMethod.ZSCORE:
        mean = np.mean(reference_values)  # E[r]
        std = np.std(reference_values)    # sigma
        if std == 0:
            score = 0.0
        else:
            score = abs(value - mean) / std  # Z = |r_obs - E[r]| / sigma
        threshold = self.zscore_threshold
    else:
        # MAD-based Z-score (robust variant)
        median = np.median(reference_values)
        mad = np.median(np.abs(reference_values - median))
        mad_adjusted = mad * self.MAD_CONSISTENCY_CONSTANT
        if mad_adjusted == 0:
            score = 0.0
        else:
            score = abs(value - median) / mad_adjusted
        threshold = self.mad_threshold
    
    is_anomaly = score > threshold
    
    return AnomalyResult(
        value=value,
        is_anomaly=is_anomaly,
        score=score,          # Z-score = surprise metric
        threshold=threshold,
        method=method,
        percentile=float(stats.percentileofscore(reference_values, value)),
    )
\end{lstlisting}
\end{implbox}

\begin{adaptbox}
The IP paper defines surprise as a Z-scored deviation: $Z = |r_{obs} - E[r]| / \sigma$. The SAP-OSS implementation directly implements this formula in \funcname{detect\_single()}, where \texttt{value} is $r_{obs}$, \texttt{mean} is $E[r]$, and \texttt{std} is $\sigma$. The computed \texttt{score} is the Z-score (surprise metric). While the IP paper applies this to memory gating, the SAP-OSS adaptation uses the same statistical framework for anomaly detection---both share the concept of identifying statistically significant deviations from expectation.
\end{adaptbox}

% =============================================================================
\section{Concept: Audit Trail with Statistical Context}
\label{sec:ip007-audit}
% =============================================================================

\begin{ipsourcebox}{IP-007, Governance Section}
\textit{``A critical aspect of innovation is the governance and auditability framework embedded within the SSG. Unlike black-box memory systems, every memory write operation generates a complete audit trail comprising the normalized deviation (Z-score), predictive model parameters at decision time, volatility estimates, and statistical significance classifications. This transforms memory from an opaque artifact into an interpretable, compliant log suitable for regulated environments.''}
\end{ipsourcebox}

\begin{tracerecord}
\textbf{IP Source ID} & IP-007 \\
\textbf{IP Location} & Governance Section, Page 6 \\
\textbf{Concept} & Audit trail with Z-score, parameters, and significance \\
\midrule
\textbf{Implementation Path} & \filepath{src/intelligence/analytics/anomaly\_detector.py} \\
\textbf{Key Files} & \filepath{anomaly\_detector.py} \\
\textbf{Functions/Classes} & \funcname{AnomalyResult}, \funcname{BatchAnomalyResult} \\
\textbf{Adaptation Type} & Structural (audit trail structure) \\
\end{tracerecord}

\begin{implbox}{src/intelligence/analytics/anomaly\_detector.py, Lines 26--46}
\begin{lstlisting}[language=Python]
@dataclass
class AnomalyResult:
    """
    Result of anomaly detection for a single value.
    
    Implements IP-007 audit trail concept:
    - score: Z-score (normalized deviation)
    - threshold: Decision threshold at detection time
    - method: Detection method used (model parameter)
    """
    value: float
    is_anomaly: bool
    score: float              # Z-score (surprise metric)
    threshold: float          # Threshold at decision time
    method: DetectionMethod   # Model parameter
    percentile: Optional[float] = None


@dataclass
class BatchAnomalyResult:
    """
    Result of batch anomaly detection.
    
    Complete audit trail per IP-007:
    - statistics: Model parameters (mean, std, etc.)
    - normality_test_pvalue: Statistical test results
    """
    values: np.ndarray
    anomaly_mask: np.ndarray
    scores: np.ndarray
    method: DetectionMethod
    normality_test_pvalue: Optional[float]  # Statistical context
    statistics: Dict[str, float]            # Full parameter snapshot
\end{lstlisting}
\end{implbox}

\begin{adaptbox}
The IP paper requires audit trails with ``Z-score, predictive model parameters at decision time, volatility estimates, and statistical significance.'' The SAP-OSS implementation captures this through the \funcname{AnomalyResult} and \funcname{BatchAnomalyResult} dataclasses, which include the Z-score (\texttt{score}), threshold at decision time (\texttt{threshold}), detection method (\texttt{method}), and complete statistics dictionary with mean, std, and Shapiro-Wilk p-value. Every detection decision is fully auditable with its statistical context.
\end{adaptbox}

% =============================================================================
\section{Concept: Confidence Calibration}
\label{sec:ip007-confidence}
% =============================================================================

\begin{ipsourcebox}{IP-007, Finance Implications}
\textit{``Confidence Calibrated by History---The memory bank's performance can be analyzed over time. The historical correlation between a high `surprise score' and a subsequent confirmed risk event informs a confidence calibration module, providing risk officers with a clear metric for when to trust an automated alert and when to escalate.''}
\end{ipsourcebox}

\begin{tracerecord}
\textbf{IP Source ID} & IP-007 \\
\textbf{IP Location} & Finance Implications, Page 4 \\
\textbf{Concept} & Historical performance tracking for confidence calibration \\
\midrule
\textbf{Implementation Path} & \filepath{src/intelligence/analytics/anomaly\_detector.py} \\
\textbf{Key Files} & \filepath{anomaly\_detector.py} \\
\textbf{Functions/Classes} & \funcname{VarianceAnalysisResult} \\
\textbf{Adaptation Type} & Structural (confidence tracking in results) \\
\end{tracerecord}

\begin{implbox}{src/intelligence/analytics/anomaly\_detector.py, Lines 366--380}
\begin{lstlisting}[language=Python]
@dataclass
class VarianceAnalysisResult:
    """
    Result of variance analysis for an account.
    
    Supports IP-007 confidence calibration through:
    - anomaly_score: Historical statistical significance
    - detection_method: Method used (for calibration grouping)
    - details: Full context for historical analysis
    """
    account_code: str
    account_name: str
    current_value: float
    prior_value: float
    variance_amount: float
    variance_percentage: float
    is_anomaly: bool
    anomaly_type: str  # spike, drop, statistical, material, new_account, none
    anomaly_score: Optional[float]     # For confidence calibration
    detection_method: Optional[str]    # For method-specific calibration
    details: Dict[str, Any]            # Full context for history
\end{lstlisting}
\end{implbox}

\begin{adaptbox}
The IP paper describes confidence calibration through historical correlation between surprise scores and confirmed events. The SAP-OSS implementation supports this through \funcname{VarianceAnalysisResult}, which captures \texttt{anomaly\_score} (the surprise metric), \texttt{detection\_method} (enabling method-specific calibration), and \texttt{details} (full context for historical analysis). This structure enables tracking of anomaly detection accuracy over time to calibrate confidence in future alerts.
\end{adaptbox}

% =============================================================================
\section{Concept: Adaptive Calibration}
\label{sec:ip007-adaptive}
% =============================================================================

\begin{ipsourcebox}{IP-007, Finance Implications}
\textit{``Continuous Model Risk Management---The SSG's internal predictive model continuously self-calibrates (\texttt{src/surprise/adaptive\_calibration.py}), adapting its parameters to changing market volatility. This process is fully logged, creating a living record of model performance and adjustment---a dynamic model risk file that demonstrates ongoing validation and governance.''}
\end{ipsourcebox}

\begin{tracerecord}
\textbf{IP Source ID} & IP-007 \\
\textbf{IP Location} & Finance Implications, Page 4 \\
\textbf{Concept} & Self-calibrating model with logged parameter evolution \\
\midrule
\textbf{Implementation Path} & \filepath{src/intelligence/analytics/anomaly\_detector.py} \\
\textbf{Key Files} & \filepath{anomaly\_detector.py} \\
\textbf{Functions/Classes} & \funcname{AnomalyDetector.\_determine\_method()} \\
\textbf{Adaptation Type} & Conceptual (adaptive method selection) \\
\end{tracerecord}

\begin{implbox}{Adaptive Method Selection}
\begin{lstlisting}[language=Python]
def _determine_method(self, values: np.ndarray):
    """
    Adaptively determine detection method based on data characteristics.
    
    Implements IP-007 adaptive calibration concept:
    - Tests distribution shape (Shapiro-Wilk)
    - Selects appropriate method (Z-score vs MAD)
    - Logs decision context (p-value) for governance
    """
    # Run Shapiro-Wilk test - adapts to current data
    _, p_value = stats.shapiro(values)
    
    # Adaptive selection based on distribution shape
    if p_value > self.shapiro_p_threshold:
        # Normal distribution - use parametric method
        return DetectionMethod.ZSCORE, p_value
    else:
        # Non-normal - adapt to robust method
        return DetectionMethod.MAD, p_value
\end{lstlisting}
\end{implbox}

\begin{adaptbox}
The IP paper describes continuous self-calibration that adapts to changing data characteristics. The SAP-OSS implementation realizes this through \funcname{\_determine\_method()}, which tests the current data distribution (Shapiro-Wilk) and adaptively selects the appropriate detection method (Z-score for normal, MAD for non-normal). While not a full online learning system, this per-batch adaptation captures the IP concept of responding to changing data characteristics rather than using fixed parameters.
\end{adaptbox}

% =============================================================================
\section{Concept: Significance Classification}
% =============================================================================

\begin{ipsourcebox}{IP-007, Hull-White Section}
\textit{``The \texttt{HullWhiteSurprise} structure includes utilities like \texttt{none()} for zeroed records and \texttt{significance\_level()} to convert Z-scores into human-readable descriptors (`insignificant' to `extreme') for logging. The \texttt{p\_value()} method uses an Abramowitz--Stegun error-function approximation to compute two-tailed probabilities, ensuring audit-ready statistics for every memory decision.''}
\end{ipsourcebox}

\begin{tracerecord}
\textbf{IP Source ID} & IP-007 \\
\textbf{IP Location} & Hull-White Section, Page 13 \\
\textbf{Concept} & Human-readable significance levels for audit \\
\midrule
\textbf{Implementation Path} & \filepath{src/intelligence/analytics/anomaly\_detector.py} \\
\textbf{Key Files} & \filepath{anomaly\_detector.py} \\
\textbf{Functions/Classes} & \funcname{VarianceAnalyzer} \\
\textbf{Adaptation Type} & Structural (anomaly type classification) \\
\end{tracerecord}

\begin{implbox}{Significance Classification}
\begin{lstlisting}[language=Python]
class VarianceAnalyzer:
    """
    Combined variance and anomaly analysis with significance classification.
    
    Implements IP-007 significance levels through anomaly_type:
    - 'spike': Significant positive change (> 25%)
    - 'drop': Significant negative change (< -25%)
    - 'statistical': Z-score/MAD triggered anomaly
    - 'material': Material balance with significant variance
    - 'new_account': New account (infinite variance)
    - 'none': No anomaly detected
    """
    
    def analyze_variance(self, ...):
        # Determine anomaly type (significance classification)
        if spike_drop_type == "spike":
            anomaly_type = "spike"
        elif spike_drop_type == "drop":
            anomaly_type = "drop"
        elif is_statistical_anomaly:
            anomaly_type = "statistical"
        elif is_material and is_material_variance:
            anomaly_type = "material"
        else:
            anomaly_type = "none"
\end{lstlisting}
\end{implbox}

\begin{adaptbox}
The IP paper describes human-readable significance levels (``insignificant'' to ``extreme'') for audit purposes. The SAP-OSS implementation adapts this through the \texttt{anomaly\_type} field in \funcname{VarianceAnalysisResult}, which classifies anomalies into business-meaningful categories: ``spike'', ``drop'', ``statistical'', ``material'', ``new\_account'', or ``none''. This provides auditors with clear, interpretable classifications rather than raw Z-scores alone.
\end{adaptbox}

% =============================================================================
\section{Implementation Summary}
% =============================================================================

\begin{table}[H]
\centering
\caption{IP-007 Concept Implementation Summary}
\label{tab:ip007-summary}
\begin{tabularx}{\textwidth}{@{}p{3cm}p{3.5cm}Xc@{}}
\toprule
\textbf{Concept} & \textbf{IP Reference} & \textbf{Implementation} & \textbf{Status} \\
\midrule
Z-score Surprise & Innovation, Page 6 & \funcname{detect\_single()} returns score & Yes \\
Audit Trail & Governance, Page 6 & \funcname{AnomalyResult} dataclass & Yes \\
Confidence Calibration & Finance, Page 4 & \funcname{VarianceAnalysisResult} & Yes \\
Adaptive Calibration & Finance, Page 4 & \funcname{\_determine\_method()} & Yes \\
Significance Levels & Hull-White, Page 13 & \texttt{anomaly\_type} field & Yes \\
\bottomrule
\end{tabularx}
\end{table}

\paragraph{Key Adaptation Pattern:}
The IP paper describes the Stochastic Surprise Gate (SSG) for \textit{neural memory systems}, using Z-scores to determine what to remember. The SAP-OSS implementation adapts this same statistical framework to \textit{anomaly detection for financial analysis}, using Z-scores to determine what is anomalous. Both share:
\begin{itemize}[nosep]
    \item Z-score as the core metric for ``significance''
    \item Statistical method selection based on distribution shape
    \item Complete audit trails with decision context
    \item Human-readable significance classifications
\end{itemize}

The IP concept of ``statistically significant deviation from expectation'' is the unifying principle across both applications.